% Encircle a target cell inline.

\newcommand{\tc}[2][black]{%
\tikz[baseline=(c.base)]{
\node[
    draw=#1,
    text=#1,
    ellipse,
    inner sep=1.5pt,
    outer sep=0pt,
    line width=0.4pt
] (c) {#2};
}%
}

\tikzset{
  stepmark/.style={
    circle, draw=black, fill=white,
    inner sep=0pt, minimum size=4.2mm,
    font=\tiny\bfseries, line width=0.4pt
  },
  steparrow/.style={
    ->,
    dashed,
    line width=0.35pt,
    >=stealth
  }
}

% ---------- Classes de couleur pour le raffinement ----------
% classRoot   : sommet 0 (distingué dès le 1er raffinement)
% classGroup  : sommets encore dans la même classe (non distingués)
% classFirst  : sommet choisi à la 1re individualisation de la branche
% classSecond : sommet choisi à la 2e individualisation de la branche
% classAuto   : sommet distingué automatiquement (dernier de sa classe)

\definecolor{classRoot}{HTML}{2166AC}
\definecolor{classGroup}{HTML}{E08214}
\definecolor{classFirst}{HTML}{B2182B}
\definecolor{classSecond}{HTML}{762A83}
\definecolor{classAuto}{HTML}{1B7837}

\tikzset{
  styGroup/.style ={draw=classGroup,  fill=classGroup!20,  text=classGroup},
  styFirst/.style ={draw=classFirst,  fill=classFirst!20,  text=classFirst,  double},
  stySecond/.style={draw=classSecond, fill=classSecond!20, text=classSecond, double},
  styAuto/.style  ={draw=classAuto,   fill=classAuto!20,   text=classAuto},
}

\newcommand{\ktgraph}[4][black]{%
\begin{tikzpicture}[
    scale=0.42,
    every node/.style={
        circle,
        draw=#1,
        fill=#1!20,
        text=#1,
        inner sep=0pt,
        minimum size=6mm,
        font=\tiny,
        line width=0.5pt
    }
]
\node (g0) at (0,0) {0};

\node[#2] (g1) at (-1.15,1.15) {1};
\node[#3] (g2) at (0,1.55) {2};
\node[#4] (g3) at (1.15,1.15) {3};

\draw[#1] (g0)--(g1);
\draw[#1] (g0)--(g2);
\draw[#1] (g0)--(g3);

\end{tikzpicture}%
}


\begin{tikzpicture}[
    every node/.style={font=\small}
]

% ---------- Initial partition ----------

\node[anchor=north,align=center] (init) at (4.75,14)
{
    \ktgraph[classGroup]{}{}{}
};

% ---------- First refinement ----------

\node[anchor=north,align=center] (root) at (4.75,10)
{
    \ktgraph[classRoot]{styGroup}{styGroup}{styGroup}
};

\draw[
    ->,
    line width=0.5pt
]
(init.south) -- (root.north);


% ---------- Level 1 ----------

\node[anchor=north,align=center] (A) at (-2.25,6)
{
    \ktgraph[classRoot]{styFirst}{styGroup}{styGroup}
};

\node[anchor=north,align=center] (B) at (4.75,6)
{
    \ktgraph[classRoot]{styGroup}{styFirst}{styGroup}
};

\node[
    anchor=north,
    align=center,
    draw=gray,
    dashed,
    rounded corners,
    inner sep=4pt
] (C) at (11.75,6)
{
    \ktgraph[classRoot]{styGroup}{styGroup}{styFirst}\\[2pt]
    {\scriptsize\itshape\color{gray} Élagage $(2\ 3)$}
};


% ---------- Leaves ----------

\node[anchor=north,align=center] (L1) at (-4,1.2)
{
    \ktgraph[classRoot]{styFirst}{stySecond}{styAuto}
};

\node[anchor=north,align=center] (L2) at (-0.5,1.2)
{
    \ktgraph[classRoot]{styFirst}{styAuto}{stySecond}
};

\node[anchor=north,align=center] (L3) at (3,1.2)
{
    \ktgraph[classRoot]{stySecond}{styFirst}{styAuto}
};

\node[
    anchor=north,
    align=center,
    draw=gray,
    dashed,
    rounded corners,
    inner sep=4pt
] (L4) at (6.5,1.2)
{
    \ktgraph[classRoot]{styAuto}{styFirst}{stySecond}\\[2pt]
    {\scriptsize\itshape\color{gray} Élagage $(1\ 3)$}
};

\node[
    anchor=north,
    align=center,
    draw=gray,
    dashed,
    rounded corners,
    inner sep=4pt
] (L5) at (10,1.2)
{
    \ktgraph[gray]{}{}{double}\\[2pt]
    {\scriptsize\itshape\color{gray} Non généré}
};

\node[
    anchor=north,
    align=center,
    draw=gray,
    dashed,
    rounded corners,
    inner sep=4pt
] (L6) at (13.5,1.2)
{
    \ktgraph[gray]{}{double}{double}\\[2pt]
    {\scriptsize\itshape\color{gray} Non généré}
};


% ---------- Edges ----------

\draw (root.south) -- (A.north)
    node[midway,above left] {$1$};

\draw (root.south) -- (B.north)
    node[midway,right] {$2$};

\draw[dashed,gray] (root.south) -- (C.north)
    node[midway,above right,gray] {$3$};


\draw (A.south) -- (L1.north)
    node[midway,left] {$2$};

\draw (A.south) -- (L2.north)
    node[midway,right] {$3$};

\draw (B.south) -- (L3.north)
    node[midway,left] {$1$};

\draw[dashed,gray] (B.south) -- (L4.north)
    node[midway,right,gray] {$3$};

\draw[dashed,gray] (C.south) -- (L5.north)
    node[midway,left,gray] {$1$};

\draw[dashed,gray] (C.south) -- (L6.north)
    node[midway,right,gray] {$2$};


% ---------- Step markers ----------

\node[stepmark] (step1)
    at ($(init.north west)+(-0.3,0.3)$) {1};

\node[stepmark] (step2)
    at ($(L1.north west)+(-0.3,0.3)$) {2};

\node[stepmark] (step3)
    at ($(L2.north east)+(-0.3,0.3)$) {3};

\node[stepmark] (step4)
    at ($(B.north east)+(-0.3,0.3)$) {4};

\node[stepmark] (step5)
    at ($(L3.north west)+(0.3,0.3)$) {5};

\node[stepmark] (step6)
    at ($(L4.north east)+(-1,0.4)$) {6};

\node[stepmark] (step7)
    at ($(C.north east)+(-1,0.4)$) {7};


% ---------- Arrows between step markers ----------

% 1 -> 2
\draw[steparrow]
(step1.south west)
to[out=225,in=135]
(step2.north west);

% 2 -> 3
\draw[steparrow]
(step2.east)
to[out=15,in=165]
(step3.west);

% 3 -> 4
\draw[steparrow]
(step3.north east)
to[out=90,in=200]
(step4.north);

% 4 -> 5
\draw[steparrow]
(step4.south east)
to[out=300,in=135]
(step5.north);

% 5 -> 6
\draw[steparrow]
(step5.east)
to[out=15,in=210]
(step6.south west);

% 6 -> 7
\draw[steparrow]
(step6.north east)
to[out=90,in=250]
(step7.south);

\end{tikzpicture}